% ============================================================
% CHAPTER 1: INTRODUCTION
% ============================================================

\chapter{Introduction}

\section{Industrial Background and Motivation}

The proliferation of electronic commerce (e-commerce) has fundamentally transformed the global retail landscape over the past two decades. According to Statista (2024), worldwide e-commerce sales surpassed \$5.7 trillion USD, accounting for approximately 20.8\% of total retail sales \cite{statista2024}. This exponential growth is driven by several converging factors: increased internet penetration, ubiquity of mobile devices, evolving consumer preferences towards convenience, and the maturation of digital payment infrastructure.

Within the e-commerce domain, the consumer electronics segment—particularly Apple products—represents a highly competitive and technically demanding niche. Apple Inc., valued at over \$3 trillion USD in market capitalization, maintains a unique position through its ecosystem integration, premium branding, and cult-like customer loyalty \cite{apple2024}. Authorized resellers and third-party retailers face significant challenges in creating digital storefronts that mirror Apple's minimalist aesthetic while delivering robust backend functionality.

From a software engineering perspective, modern e-commerce platforms are no longer monolithic web applications but rather complex distributed systems requiring:

\begin{itemize}
    \item \textbf{High concurrency handling:} Supporting thousands of simultaneous users during product launches or promotional events (e.g., Black Friday sales).
    \item \textbf{Data consistency:} Ensuring ACID compliance for inventory management and financial transactions to prevent overselling or payment discrepancies.
    \item \textbf{Security robustness:} Protecting Personally Identifiable Information (PII) and payment credentials against sophisticated attack vectors (SQL injections, XSS, CSRF).
   

 \item \textbf{User experience optimization:} Achieving sub-2-second page load times and intuitive interfaces to minimize cart abandonment (industry average: 69.57\% \cite{baymard2023}).
\end{itemize}

The confluence of these requirements positions e-commerce system development as an ideal case study for applying modern software architecture principles, secure coding practices, and full-stack development methodologies—all core competencies expected of undergraduate Computer Science graduates entering the industry.

\section{Problem Statement}

Despite the maturity of e-commerce as a domain, many existing platforms—particularly those developed by small-to-medium enterprises (SMEs) or used in academic contexts—suffer from critical deficiencies that hinder scalability, security, and user satisfaction. Through preliminary market analysis and technical audits of comparable student projects, the following recurring issues were identified:

\subsection{Architectural and Performance Limitations}

\begin{enumerate}
    \item \textbf{Monolithic Bottlenecks:} Legacy architectures tightly couple frontend presentation, business logic, and data access layers into a single codebase. This design pattern, while simple for initial development, creates scaling bottlenecks. For instance, a surge in user traffic (e.g., due to viral marketing) can cause the entire application to fail, rather than isolating load to specific microservices.
    
    \item \textbf{Inefficient Database Queries:} Lack of proper indexing strategies (e.g., B-Tree indexes on frequently queried columns like product names or categories) results in query execution times exceeding 500ms under moderate load, degrading user experience. 
    
    \item \textbf{Absence of Caching Mechanisms:} Repeated database hits for static data (product catalogs, category trees) without implementing Redis or in-memory caching unnecessarily burdens the database server.
\end{enumerate}

\subsection{Security Vulnerabilities}

\begin{enumerate}
    \item \textbf{Weak Authentication Protocols:} Many systems rely on session-based authentication with cookies vulnerable to session hijacking and Cross-Site Request Forgery (CSRF) attacks. Modern token-based authentication (e.g., JSON Web Tokens - JWT) with short expiration windows and refresh token rotation is rarely implemented correctly.
    
    \item \textbf{Inadequate Input Validation:} Failure to sanitize user inputs exposes systems to injection attacks (SQL injection via search queries,XSS via product reviews). The OWASP Top 10 consistently lists injection flaws as the #1 web application vulnerability \cite{owasp2021}.
    
    \item \textbf{Plaintext Password Storage:} Surprisingly, a non-trivial number of educational projects store passwords in plaintext or use weak hashing algorithms (e.g., MD5), violating basic security hygiene.
\end{enumerate}

\subsection{User Experience Deficiencies}

\begin{enumerate}
    \item \textbf{Non-Responsive Design:} Failure to implement mobile-first or adaptive layouts alienates the 59\% of global e-commerce traffic originating from mobile devices \cite{statcounter2024}.
    
    \item \textbf{Poor Checkout Flow:} Complicated multi-step checkout processes with mandatory account creation increase cart abandonment. Research indicates that 24\% of users abandon carts due to being forced to create an account \cite{baymard2023}.
    
    \item \textbf{Lack of Intelligent Features:} Absence of AI-driven search, chatbot assistance, or personalized recommendations reduces competitive advantage against market leaders (Amazon, Shopify).
\end{enumerate}

\subsection{Integration and Maintainability Issues}

\begin{enumerate}
    \item \textbf{Tight Coupling with Third-Party Services:} Hardcoding payment gateway or logistics APIs without abstraction layers (Strategy/Adapter design patterns) makes it difficult to swap providers or handle API deprecations.
    
    \item \textbf{Insufficient Documentation:} Lack of inline comments, API documentation (e.g., Swagger/OpenAPI), and architectural diagrams hampers onboarding of new developers and future maintainability.
\end{enumerate}

\vspace{0.5cm}
\noindent Addressing these deficiencies through a meticulously designed, industry-aligned e-commerce system forms the core motivation for this undergraduate thesis project.

\section{Research Objectives}

The overarching goal of this project is to design, implement, and rigorously evaluate a full-stack web-based e-commerce platform tailored for Apple product retail, adhering to industrial best practices and modern software engineering standards. The system shall serve as both a functional prototype and an academic artifact demonstrating the application of theoretical knowledge to real-world problem domains.

\subsection{Primary Objectives}

\begin{enumerate}
    \item \textbf{To architect a scalable and secure backend infrastructure} utilizing **Node.js** and **Express.js** framework, implementing RESTful API design principles with JWT-based stateless authentication to enable horizontal scalability.
    
    \item \textbf{To design and implement a robust database schema} using **MongoDB** (NoSQL) with proper indexing strategies, ensuring data integrity through validation middleware and Mongoose schema constraints. The schema shall normalize data to at least **Third Normal Form (3NF)** equivalent where applicable to relational concepts.
    
    \item \textbf{To develop a modern, responsive frontend} using semantic **HTML5**, modular **CSS3** (with Tailwind utility framework), and vanilla **JavaScript** (ES6+), ensuring cross-browser compatibility and mobile-first design philosophy.
    
    \item \textbf{To integrate advanced features} including:
    \begin{itemize}
        \item Google OAuth 2.0 for social authentication
        \item AI-powered chatbot using Google Gemini API for customer support
        \item 3D product visualization using Three.js WebGL library
        \item Gamified loyalty program with points and rank system
    \end{itemize}
    
    \item \textbf{To implement a comprehensive admin dashboard} with real-time analytics (Chart.js), inventory management with stock tracking, order lifecycle management, and voucher administration.
    
    \item \textbf{To rigorously test system functionality, security, and performance}, documenting results through:
    \begin{itemize}
        \item Load testing with Apache JMeter (simulating 500+ concurrent users)
        \item Security audits using OWASP ZAP automated scanning
        \item User Acceptance Testing (UAT) with System Usability Scale (SUS) evaluation
    \end{itemize}
\end{enumerate}

\subsection{Academic Contributions}

Beyond functional deliverables, this thesis aims to contribute:

\begin{enumerate}
    \item \textbf{A reference architecture} for junior developers and students learning full-stack MERN-like development (MongoDB, Express, React concepts via vanilla JS, Node.js).
    
    \item \textbf{Comparative analysis} of architectural patterns (Monolithic vs. Modular Monolith vs. Microservices) in the context of SME e-commerce requirements.
    
    \item \textbf{Cost-benefit analysis} comparing custom development against SaaS platforms (Shopify, WooCommerce) to inform strategic IT decisions for startups.
\end{enumerate}

\section{Scope and Delimitations}

To maintain focus and feasibility within the constraints of an undergraduate thesis timeline and resources, the project scope is explicitly defined below.

\subsection{Inclusions (In Scope)}

\begin{itemize}
    \item \textbf{Frontend Components:}
    \begin{itemize}
        \item Landing page with 3D product showcase (Three.js)
        \item Product listing with advanced filtering (category, price range) and sorting
        \item Shopping cart with persistent storage (LocalStorage + server sync)
        \item Checkout workflow with guest and authenticated user support
        \item User account dashboard (order history, profile management, voucher redemption)
        \item Admin dashboard with analytics and CRUD operations
    \end{itemize}
    
    \item \textbf{Backend Features:}
    \begin{itemize}
        \item User registration and authentication (JWT + Google OAuth 2.0)
        \item Role-Based Access Control (RBAC): User vs. Admin
        \item Product catalog management with text search (MongoDB text indexes)
        \item Order processing with voucher application and inventory deduction
        \item Gamified loyalty system (points accrual, rank progression: Silver → Gold → VIP)
        \item AI chatbot integration (Google Gemini API) for product queries
    \end{itemize}
    
    \item \textbf{Database Design:}
    \begin{itemize}
        \item NoSQL schema design (MongoDB) with embedded documents and references
        \item Collections: Users, Products, Orders, Carts, Vouchers
        \item Text indexes for search optimization
        \item Timestamps and audit fields (createdAt, updatedAt)
    \end{itemize}
\end{itemize}

\subsection{Exclusions (Out of Scope)}

\begin{itemize}
    \item \textbf{Real Payment Processing:} Actual financial transactions with banks or payment gateways (Stripe, PayPal) are **simulated**. Only Cash-on-Delivery (COD) is implemented as a payment method. This limitation is justified by:
    \begin{itemize}
        \item Regulatory compliance requirements (PCI-DSS certification)
        \item Legal liability concerns for handling real financial data
        \item Requirement for business registration and merchant accounts
    \end{itemize}
    
    \item \textbf{Logistics Integration:} Real-time shipment tracking via carrier APIs (FedEx, DHL, Vietnam Post) is **not integrated**. Order status updates are manual (admin dashboard). Rationale: API costs and complexity beyond academic scope.
    
    \item \textbf{Infrastructure Deployment:} While the system is **designed** for cloud deployment (Docker-ready architecture), production deployment on AWS/GCP/Azure with auto-scaling, load balancers, and CDN configuration is **not executed** due to cost constraints. The system is demonstrated on **localhost** or a single-node **Render.com** free-tier instance.
    
    \item \textbf{Internationalization (i18n):} Multi-language support and currency conversion are **not implemented**. The system targets Vietnamese market (VND currency, Vietnamese language for some UI elements).
    
    \item \textbf{Advanced AI Features:} While a chatbot is integrated, advanced recommendation systems (collaborative filtering, deep learning models) are **out of scope**.
\end{itemize}

\section{Development Methodology}

This project adopts an **Agile development methodology**, specifically leveraging concepts from **Scrum** to facilitate iterative development, continuous feedback, and adaptive planning. While a formal Scrum team structure (Product Owner, Scrum Master, Development Team) is not applicable to a solo academic project, the core principles are adapted as follows:

\subsection{Sprint Structure}

Development is divided into **weekly sprints**, each with defined goals and deliverables:

\begin{itemize}
    \item \textbf{Sprint 1-2:} Requirement analysis, technology evaluation, database schema design, project setup
    \item \textbf{Sprint 3-4:} Backend API development (authentication, product CRUD)
    \item \textbf{Sprint 5-6:} Frontend development (landing page, product listing)
    \item \textbf{Sprint 7-8:} Cart and checkout implementation
    \item \textbf{Sprint 9-10:} Admin dashboard development
    \item \textbf{Sprint 11-12:} Testing, debugging, documentation, thesis writing
\end{itemize}

\subsection{Development Workflow (Based on Git Flow)}

\vspace{0.3cm}
\begin{lstlisting}[language=bash, caption=Git Branching Strategy Example]
# Feature development
git checkout -b feature/user-authentication
# ... Development work ...
git commit -m "Implement JWT middleware"
git push origin feature/user-authentication

# Merge to main after testing
git checkout main
git merge feature/user-authentication
\end{lstlisting}

Version control is managed via **GitHub** with the following branch strategy:

\begin{itemize}
    \item \texttt{main}: Production-ready code
    \item \texttt{feature/*}: Individual feature development (e.g., \texttt{feature/chatbot-integration})
    \item Commits follow **Conventional Commits** format (\texttt{feat:}, \texttt{fix:}, \texttt{docs:})
\end{itemize}

\subsection{Tools and Infrastructure}

\begin{table}[H]
\centering
\caption{Development Tools and Technologies}
\label{tab:dev_tools}
\begin{tabular}{@{}p{0.25\linewidth}p{0.30\linewidth}p{0.35\linewidth}@{}}
\toprule
\textbf{Category} & \textbf{Tool/Technology} & \textbf{Purpose} \\
\midrule
IDE & Visual Studio Code & Code editing with ESLint, Prettier \\
Version Control & Git + GitHub & Source code management \\
API Testing & Postman & Endpoint testing and documentation \\
Database GUI & MongoDB Compass & Database inspection and queries \\
Design & Figma & UI/UX wireframing (optional) \\
Diagramming & draw.io, PlantUML & UML diagrams, ERD \\
Documentation & LaTeX (Overleaf) & Thesis report generation \\
Deployment & Render.com (Free tier) & Live demo hosting \\
\bottomrule
\end{tabular}
\end{table}

\subsection{Quality Assurance Practices}

\begin{itemize}
    \item \textbf{Code Review:} Self-review via GitHub PR interface before merging features
    \item \textbf{Linting:} ESLint configured with Airbnb style guide
    \item \textbf{Testing:} Manual testing supplemented with Postman automated test collections
    \item \textbf{Security Scanning:} Periodic npm audit and OWASP ZAP scans
\end{itemize}

\section{Expected Outcomes and Deliverables}

Upon successful completion, the project will yield the following tangible and intangible deliverables:

\subsection{Functional Deliverables}

\begin{enumerate}
    \item \textbf{Fully Functional Web Application:} Accessible via URL (e.g., \texttt{https://project3-icy1.onrender.com}), demonstrating all in-scope features with sample product data (20+ products across 5 categories).
    
    \item \textbf{Source Code Repository:} Well-documented GitHub repository with:
    \begin{itemize}
        \item README with setup instructions
        \item Code comments and JSDoc annotations
        \item \texttt{package.json} with dependency specifications
        \item Environment variable template (\texttt{.env.example})
    \end{itemize}
    
    \item \textbf{Database Dump:} MongoDB export file (\texttt{mongodump}) containing sample data for reproducibility.
\end{enumerate}

\subsection{Technical Documentation}

\begin{enumerate}
    \item \textbf{System Architecture Document:} Detailed diagrams (deployment, component, sequence) explaining system design.
    
    \item \textbf{API Specification:} Postman collection or OpenAPI (Swagger) YAML file documenting all endpoints.
    
    \item \textbf{Database Schema Documentation:} ERD with field descriptions, data types, and relationships.
    
    \item \textbf{User Manual:} Step-by-step guide for end users (customers and admins).
\end{enumerate}

\subsection{Academic Deliverables}

\begin{enumerate}
    \item \textbf{Thesis Report:} This comprehensive document (60+ pages) detailing every phase of the software development lifecycle.
    
    \item \textbf{Presentation Slides:} PowerPoint/PDF deck for thesis defense (20-25 slides).
    
    \item \textbf{Demo Video:} 10-minute screencast demonstrating core features.
\end{enumerate}

\section{Report Organization}

This thesis is structured into seven chapters, methodically documenting the entire project lifecycle:

\begin{itemize}
    \item \textbf{Chapter 1: Introduction} (current chapter) establishes the industrial context, problem domain, objectives, scope, and methodology.
    
    \item \textbf{Chapter 2: Technology Overview} provides a comprehensive survey of the technology stack (MERN-like ecosystem), including comparative analysis justifying technology choices.
    
    \item \textbf{Chapter 3: System Analysis} details the requirements engineering process, presenting functional and non-functional requirements, use cases, and business workflows.
    
    \item \textbf{Chapter 4: System Design and Architecture} presents the technical blueprint: architecture diagrams (C4 model), database schema (ERD), API design, and security architecture.
    
    \item \textbf{Chapter 5: Implementation and Evaluation} describes the coding phase, presenting key code snippets (with explanations), followed by rigorous testing results (performance benchmarks, security audits, user feedback).
    
    \item \textbf{Chapter 6: Discussion} critically analyzes the project outcomes, comparing against market solutions, discussing technical trade-offs, and identifying limitations.
    
    \item \textbf{Chapter 7: Conclusion and Future Work} summarizes achievements, academic contributions, and proposes a roadmap for future enhancements.
    
    \item \textbf{Appendices} include selected source code listings, API endpoint tables, and test case specifications.
\end{itemize}

Each chapter builds upon the previous, forming a cohesive narrative that mirrors the actual development process while meeting academic rigor standards for undergraduate thesis work in Computer Science.

\vspace{1cm}
\noindent \textbf{Note:} All technical claims, performance metrics, and architectural decisions documented in this report are based on the actual implemented system, with code citations provided where applicable to ensure traceability and academic integrity.

% End of Chapter 1
